\subsection{K-Prototypes Pipeline Implementation}

The clustering pipeline consists of a script file called \texttt{k\_prototypes.py}. This script can be executed using \texttt{python k\_prototypes.py} in the appropriate directory and is not connected to FastAPI since it is time-consuming and should be done offline. The function for clustering is located in the \texttt{kmodes} at \texttt{kmodes.kprototypes.KPrototypes}. Everything else (feature extraction, scaling, projection, evaluating, plotting, persisting, and summarising with AI) is part of the wrapper around this algorithm. The entire wrapper is part of the \texttt{ml-batch} microservice.

\subsubsection*{Feature Construction}
\label{ch:feature_construction}

The function \texttt{build\_patient\_features(engine)} invokes an array of standalone SQL queries using \texttt{pandas.read\_sql(text(\dots), engine)}. Each query retrieves a compact dataframe, indexed by \texttt{patient\_id}, for each distinct feature category (patient demographics, encounters, condition flags, average BMI/BP/glucose, care plan, physicians, off-hours flag, substance abuse, pain level, readmission level, and so forth). The function performs a chain of \texttt{df.merge(right\_df, on="patient\_id", how="left")} operations on those dataframes, and then applies type-aware imputation.

Here is an example of the SQL magic, illustrating the process used to derive the \emph{dominant care plan} categorical variable (which constitutes one of the cluster-determining features). The 40 care plan categories are bucketized into 7 clinical categories based on severity rank with \texttt{ROW\_NUMBER() OVER (PARTITION BY patient\_id ORDER BY severity\_rank ASC)}.

\begin{lstlisting}[style=sqlstyle]
WITH CategorizedPlans AS (
    SELECT patient_id,
           CASE
               WHEN description IN (
                   'Cancer care plan (record artifact)',
                   'Dialysis care plan (record artifact)',
                   'Terminal care (regime/therapy)'
               ) THEN 'Complex, Oncology & Palliative'
               WHEN description = 'Routine antenatal care (regime/therapy)'
                                 THEN 'Maternal & Antenatal'
               WHEN LOWER(description) ~ '(fracture|wound|surgery|injury|trauma)'
                                 THEN 'Acute Trauma, Wound & Surgery'
               WHEN LOWER(description) ~ '(respiratory|asthma|pulmonary|covid)'
                                 THEN 'Respiratory & Infectious'
               WHEN LOWER(description) ~ '(diabetes|heart|hypertension|cardiac)'
                                 THEN 'Chronic Cardio-Metabolic'
               ELSE 'General Rehab & Routine'
           END AS care_plan_category,
           CASE
               WHEN description IN (
                   'Cancer care plan (record artifact)',
                   'Dialysis care plan (record artifact)',
                   'Terminal care (regime/therapy)'
               ) THEN 1
               WHEN description = 'Routine antenatal care (regime/therapy)' THEN 2
               WHEN LOWER(description) ~ '(fracture|wound|surgery|injury)'   THEN 3
               WHEN LOWER(description) ~ '(respiratory|asthma|pulmonary)'    THEN 4
               WHEN LOWER(description) ~ '(diabetes|heart|hypertension)'     THEN 5
               ELSE 6
           END AS severity_rank
    FROM careplan
),
RankedPatientPlans AS (
    SELECT patient_id, care_plan_category,
           ROW_NUMBER() OVER (
               PARTITION BY patient_id ORDER BY severity_rank ASC
           ) AS rank_id
    FROM CategorizedPlans
)
SELECT patient_id, care_plan_category AS dominant_careplan
FROM RankedPatientPlans
WHERE rank_id = 1;
\end{lstlisting}

The same technique is used for the other derived categories as well: \texttt{substance\_use\_profile}, \texttt{pain\_management\_tier}, \texttt{readmission\_risk\_tier}, \texttt{cholesterol\_risk\_category}, \texttt{respiratory\_disease\_profile}, and many more.

\subsubsection*{Feature Bucketing}

After the feature frames are merged, the preprocessing steps utilize an intentional approach for bucketing and imputation for dealing with the sparsity of clinical features:

\begin{lstlisting}[style=pythonstyle]
for col in ["wellness_count", "inpatient_count", ...]:
    df[col] = df[col].fillna(0)        # zero is semantically correct
df["dominant_careplan"] = df["dominant_careplan"].fillna("None")
df["bmi_type"] = pd.cut(df["avg_bmi"],
                        bins=[-inf, 18.5, 25, 30, 35, 40, inf],
                        labels=["Underweight", "Normal", "Overweight",
                                "Obese I", "Obese II", "Obese III"]) \
                  .astype(object).where(df["avg_bmi"].notna(), "Unknown")
\end{lstlisting}

\noindent Numerical features that are sparse, such as body mass index, are intentionally bucketed with \texttt{pd.cut} rather than being filled with their mean value. Given that K-Prototypes assumes that numerical centroids represent real-world measurements, adding the mean artificially would result in clusters being affected. These are then grouped into clinical buckets and any missing information is sent to a \texttt{"Unknown"} group. On the other hand, count-type features can be simply imputed with zeros since this accurately represents a lack of events.

\subsubsection*{The Role of \texttt{log1p}}

\texttt{total\_\allowbreak medications}, \texttt{lifetime\_\allowbreak medication\_\allowbreak days}, and \texttt{total\_\allowbreak encounter\_\allowbreak duration\_\allowbreak hours} have strong right tails: a few cases with many medications or long durations have values an order of magnitude above the median. For those columns, there are two steps in the pipeline where \texttt{np.log1p} is used in order to reduce such extreme values:

\begin{lstlisting}[style=pythonstyle]
# inside remove_outliers_zscore; z-scores on raw skewed counts would
# flag a quarter of the cohort as outliers purely on distribution shape.
for i, col in enumerate(cols):
    if col in _SKEWED_OUTLIER_COLS:
        data[:, i] = np.log1p(np.clip(data[:, i], 0, None))

# inside run_kprototypes, so the scaler does not push the "normal" mass
# of patients into a tiny corner of [0, 1].
for col in ["total_medications", "lifetime_medication_days",
            "total_encounter_duration_hours"]:
    if col in numerical_cols:
        df_scaled[col] = np.log1p(df_scaled[col].clip(lower=0).astype(float))
\end{lstlisting}

\noindent \texttt{log1p} (\(\log(1+x)\)) is better than simply taking \(\log(x)\) because \texttt{log1p} is defined for \(x=0\) (and the mode of these features is zero). This transformation is done only on \texttt{df\_scaled} (the input matrix for the K-Prototypes algorithm) and \emph{not} on \texttt{df\_work} (the frame returned for the LLM prompt and the database).

\subsubsection*{PCA for Severity Compression, FAMD for Visualisation}

There are two distinct dimensionality reduction methods applied for two very different applications. In broad terms, the dimensionality reduction methods, PCA, and FAMD perform two distinct roles: Feature Compression and Visualization. On a fundamental level, the methods of feature compression transform an original set of features that describe a patient through a mathematical process to represent him/her in a low-dimensional space. With careful consideration to maintain maximum variability and correlation between the variables, this technique reduces an extremely complicated set of features describing a patient to a single $(x,y)$ coordinate pair on a 2D plane.

\vspace{0.5em}
\textbf{PCA} (principal component analysis) accomplishes this task for continuous variables by determining an orthogonal basis for the input space that orders axes based on their variance contribution. In this case, applied to the three correlated clinical severity scores, the first principal component contains the common "how sick is this patient" information in one variable that can be weighted against other numericals by the \textbf{K-Prototypes} algorithm. As a safeguard, configurations containing \texttt{clinical\_severity\_score} and any of its components are disallowed, as this would be redundant variance.

\paragraph{Principal component analysis for the three components representing severity.} If the \texttt{clinical\_severity\_score} composite feature is present in the configuration, it is obtained by applying the \texttt{log1p} function to the three raw severity features (\texttt{total\_procedures}, \texttt{total\_medications}, \texttt{active\_conditions\_count}), followed by standardisation, and then projecting them onto the first principal component. Since the resulting axis is already a z-score by construction, the outer min-max scaler skips it:

\begin{lstlisting}[style=pythonstyle]
sev_input        = np.log1p(df[sev_cols].clip(lower=0).values.astype(float))
sev_standardized = StandardScaler().fit_transform(sev_input)
sev_pca          = PCA(n_components=1)
css              = sev_pca.fit_transform(sev_standardized)[:, 0]
print(f"explained variance ratio: {sev_pca.explained_variance_ratio_[0]:.4f}")

cols_to_scale = [c for c in numerical_cols if c != "clinical_severity_score"]
df_scaled[cols_to_scale] = scaler.fit_transform(df_scaled[cols_to_scale].values)
\end{lstlisting}

\paragraph{FAMD for the 2-D visualisation.} The scatter plot used in the dashboard is not computed using PCA but rather \textbf{FAMD} (Factor Analysis of Mixed Data), which comes from the \texttt{prince} library. This approach is simply the natural extension of PCA to the case where there is a mixture of quantitative and qualitative data types; quantitative data is standardised, while qualitative data is one-hot-encoded and then reweighted based on their inverse category frequencies.

\begin{lstlisting}[style=pythonstyle]
import prince
famd   = prince.FAMD(n_components=2, random_state=42).fit(df_feat)
coords = famd.transform(df_feat).values  # shape (n, 2)
\end{lstlisting}

PCA can be applied to numerics only; applying an ordinal encoding to categoricals prior to PCA destroys the categorical distance geometry utilized by K-Prototypes in its clustering process. By giving both types of variables equal treatment, FAMD creates a projection that is faithful to the space being clustered into.

\subsubsection*{Scaling the Numerical Block}

Three numerical scalers are exposed through a name-to-instance map and chosen per run. Each one is taken straight from \texttt{sklearn.preprocessing} and rewrites the column with a different rule:

\begin{lstlisting}[style=pythonstyle]
_SCALER_MAP = {
    "minmax":   MinMaxScaler(),    # x' = (x - min) / (max - min) -> [0, 1]
    "robust":   RobustScaler(),    # x' = (x - median) / IQR      -> centered on 0, unbounded
    "standard": StandardScaler(),  # x' = (x - mean)   / std      -> ~ N(0, 1), unbounded
}
\end{lstlisting}

\paragraph{MinMaxScaler.} Squeezes each column linearly into the finite interval \([0, 1]\). The minimum value maps to zero, the maximum to one, and all other values map linearly within that range. The shape of the distribution is preserved exactly, with only the support changing.

\paragraph{StandardScaler.} Centers each column around its own mean and divides by its own standard deviation, resulting in a column whose approximate distribution is \(\mathcal{N}(0, 1)\). The centering and unit variance are quite useful in practice, but the unbounded nature of the Gaussian causes the heavy tails in the column to remain unaltered.

\paragraph{RobustScaler.} Centers around the median and divides by the inter-quartile range (IQR). Both of these measures are robust to outliers, making this the best scaler to use when there are any real outliers present in the data; however it is still an unbounded output.

\paragraph{Why MinMax works best for K-Prototypes.} The K-Prototypes objective consists of a squared-Euclidean distance between the numerical features and their corresponding cluster prototypes, plus a \(\gamma\)-scaled Hamming distance between the categorical features and their corresponding cluster prototypes, where the Hamming term is bounded in \([0, n_{\text{cat}}]\). By scaling the numerical features to the same \([0, 1]\) interval using \texttt{MinMaxScaler}, the two terms have similar magnitudes no matter how far apart the numbers are from each other, hence \(\gamma\) can be automatically tuned by \texttt{kmodes} without any one term overpowering the other. Although \texttt{StandardScaler} centers at zero, it does not eliminate the long tails in the distribution, which increases the importance of outlier-heavy features for the numerical objective. \texttt{RobustScaler}, while more robust to outliers, still has a few outlying samples with extremely high or low numerical values (not within \([0, 1]\)), which resulted in less stable clusters on average across the benchmark.

\subsubsection*{Z-score Outlier Detection}
\label{ch: z_score}

Before the process of clustering begins, function \texttt{remove\_outliers\_zscore()} scans each patient according to a static list of numerical attributes. Z-Score of any value \(x\) means how far away \(x\) is from the average of the respective column measured by standard deviations. Thus, a patient with \(z=6\) on the \texttt{lifetime\_medication\_days} feature will be six standard deviations greater than the average value in this column. Outliers' scan goes over the original numerical \texttt{numpy} array of the patient data frame (NOT over the output of cluster analysis - since there are no clusters formed yet), but with an initial \texttt{log1p()} transformation applied to those columns with strong right skewness.

\begin{lstlisting}[style=pythonstyle]
data = df[cols].to_numpy(dtype=float)
for i, col in enumerate(cols):
    if col in _SKEWED_OUTLIER_COLS:
        data[:, i] = np.log1p(np.clip(data[:, i], 0, None))

mask = np.ones(len(df), dtype=bool)
for i, col in enumerate(cols):
    z     = np.abs((data[:, i] - np.nanmean(data[:, i])) / np.nanstd(data[:, i]))
    mask &= z < threshold      # threshold = 6 in production
\end{lstlisting}

The patient is retained iff \emph{all} columns' z-scores remain under the threshold. Notably, outliers are not removed, because the function produces two frames, and the \texttt{df\_outliers} one is propagated throughout the processing chain and receives the label value \(-1\) upon persistence. The result is that the outliers form a distinct cluster on the dashboard (\emph{outliers}), are explained by the LLM using the same prompt as a normal cluster, and even mined separately as a unique scope by PrefixSpan. No patient is thereby ignored in analysis; at the same time, the centroid for each of the four actual clusters represents an ordinary patient, not outliers.

\subsubsection*{Running the Algorithm}

The clustering call will build a \texttt{numpy} matrix where numeric columns appear before string columns; the latter will be stored as strings and their column indices will be passed through \texttt{categorical=}. It’s important to note that constructing the matrix piece by piece ensures that the Cython loop inside \texttt{kmodes} runs quickly:

\begin{lstlisting}[style=pythonstyle]
num_matrix    = df_scaled[numerical_cols].to_numpy(dtype=np.float64)
cat_matrix    = df_scaled[categorical_cols].astype(str).to_numpy()
matrix        = np.concatenate([num_matrix, cat_matrix], axis=1)
cat_positions = list(range(len(numerical_cols),
                           len(numerical_cols) + len(categorical_cols)))
kproto = KPrototypes(n_clusters=k, init="Huang",
                     random_state=rs, n_jobs=-1, gamma=gamma)
clusters = kproto.fit_predict(matrix, categorical=cat_positions)
\end{lstlisting}

\texttt{init="Huang"} is the frequency-based initialization scheme from Huang (1998), which seeds prototypes using attribute frequency distributions, and combined with a fixed \texttt{random\_state} ensures reproducibility across runs. If \texttt{gamma=None}, \texttt{kmodes} computes \(\gamma\) by using the standard deviation of the (previously min-max normalized) numeric portion, ensuring that both parts of the objective function are balanced.

\subsubsection*{Configuration Loop and Persistence to the Database}

A list of feature configurations is stored in \texttt{configs.py} using the format \texttt{CONFIGS = [\{"version": \dots, \allowbreak "preferred\_\allowbreak k": k, \allowbreak "numerical": [\dots], \allowbreak "categorical": [\dots]\}]}. \textbf{63 configurations} (\texttt{v1} to \texttt{v63}) have been created during the experimentation process, where each configuration tests a new hypothesis (such as severity-only, utilization-only, behaviour-focused, cardio-metabolic, demographic baseline only, and many more - the file archives all retired configurations with comments for an audit trail of the search process). The complete benchmarking for all 63 configurations and the reasoning behind choosing \texttt{v55} as the champion model is explained in Chapter~\ref{ch:testing}; the selected configuration remains uncommented in the code:

\begin{lstlisting}[style=pythonstyle]
{
    "version": 55,
    "preferred_k": 4,
    "numerical": [
        "age",
        "active_conditions_count",
        "total_medications",
        "total_procedures",
        "preventive_to_reactive_ratio",
    ],
    "categorical": [
        "gender",
        "pregnancy_history",
        "has_heart_disease",
        "has_diabetes",
        "has_hypertension",
        "condition_burden",
        "bmi_type",
        "dominant_careplan",
    ],
},
\end{lstlisting}

The primary loop of \texttt{k\_prototypes.py} runs over the Cartesian product of \texttt{THRESHOLDS}, \texttt{CONFIGS}, \texttt{SCALERS}, \texttt{GAMMAS}, \texttt{RANDOM\_STATES}, and the selected \(k\)s. Even when the production configuration forces all axes to have only one value each, this is the construct that makes parameter tuning possible, simply replace any of those singletons with a non-singleton list and you will produce a benchmark sweep:

\begin{lstlisting}[style=pythonstyle]
for threshold in THRESHOLDS:                          # outlier sensitivity
    df_cleaned, df_outliers = remove_outliers_zscore(df, threshold=threshold)
    for cfg in CONFIGS:                               # 63 versions
        for scaler_name in SCALERS:                   # minmax / robust / standard
            for gamma in GAMMAS:                      # auto or fixed
                for rs in RANDOM_STATES:              # init seeds
                    for k in k_values:                # cluster count
                        df_result, X_num_scaled, clusters, _, _ = run_kprototypes(...)
                        report, eval_metrics = full_evaluation_report(...)
                        # one row per run, written incrementally
                        general_rows.append({"version": ..., **eval_metrics})
                        flush_general_conclusion()
\end{lstlisting}

\noindent Each iteration produces its own \texttt{conclusions.txt}, \texttt{conclusions.xlsx}, FAMD scatter plot, heatmap, bar plot, and stacked categorical distribution inside a folder named for the particular iteration inside \texttt{outputs/}; however, a unique \texttt{general\_conclusion.xlsx} combines all the metrics from each run in order to make comparison easier. Most importantly, all metrics used for evaluation (silhouette, Gower silhouette, Calinski--Harabasz, and Davies--Bouldin, both in raw numerical space and in the FAMD projection) are recalculated each time with \texttt{full\_evaluation\_report}, which ensures that there is a complete signature for each run in the benchmark Excel file. The detailed result analysis and selection of \texttt{v55} are presented in Chapter~\ref{ch:testing}.

\subsubsection*{From DataFrame to \texttt{numpy} and Back to Words}

The K-Prototypes algorithm and the sklearn metrics use \texttt{numpy} arrays, but the prompts from the LLM model, the cards from the dashboard, and the \texttt{cluster\_\allowbreak patient.features} JSONB field require the original strings to be restored (\texttt{gender = "F"}, \texttt{bmi\_\allowbreak type = "Obese II"}, \texttt{dominant\_\allowbreak careplan = "Chronic Cardio-Metabolic"}). This is achieved by having \texttt{run\_\allowbreak kprototypes} maintain two frames simultaneously, one \emph{raw} (\texttt{df\_\allowbreak work}), which is returned to the persistence layer and LLM prompt, and the other (\texttt{df\_\allowbreak scaled}) containing only the transformations required for clustering (log1p, severity PCA, scaler), which is converted to the \texttt{numpy} array. The integer array of cluster labels from the \texttt{fit\_\allowbreak predict} call is then appended to the raw frame as a single column; no further "decoding" of the results is necessary.

\begin{lstlisting}[style=pythonstyle]
df_work   = df.copy()        # raw values, go to the LLM and the DB
df_scaled = df_work.copy()   # transformed copy, goes to numpy / sklearn

# ... log1p, severity PCA, scaler.fit_transform applied to df_scaled only ...

clusters = kproto.fit_predict(matrix, categorical=cat_positions)  # ndarray[int]
df_work["cluster"] = clusters    # one extra column, original strings intact
\end{lstlisting}


\subsubsection*{Reuse of Derived Features Across Configurations}

The 21 SQL features generated by \texttt{build\_patient\_features} are calculated \emph{only once} for each run and stored within the \texttt{df} dataframe. In the configuration loop, we simply \emph{select} the columns we need from that dataframe; no SQL will be executed when going from one configuration \(v_i\) to the next \(v_{i+1}\). That's why 63-config sweeping is possible - the computation of heavy lifting (extracting the rows, LEFT JOIN, bucketing categoricals) is done once, and then K-Prototypes fitting becomes an in-memory numpy calculation on the very same frame.

The same logic is employed on the persistence end: the \texttt{save\_cluster\_results} function is passed \texttt{extra\_feature\_cols=[...]}, so all the columns that aren't used in the clustering (like \texttt{avg\_bmi}, \texttt{qaly}, \texttt{has\_cancer}, and \texttt{smoking\_status}) are saved within the \texttt{cluster\_patient.features} JSONB column. The cluster detail and patient detail views in the frontend read the JSONB, which means the non-used-by-clustering features will be saved for use with the analytical API and summaries for each AI prompt.

\subsubsection*{AI-driven Cluster Cards on the Frontend}

Once the clustering step is done, \texttt{get\_clusters\_summary} generates a prompt that combines information about each cluster’s numeric mean, categorical composition, top conditions, medications, procedures, care plan distribution, and population context into one string and sends it via \texttt{ask\_ai} (same OpenAI-powered function the SQL agent uses). The prompt requests the following from the model:

\begin{lstlisting}[style=pythonstyle]
[Cluster ID] | [#HexColor] | [Short Title] | [1-2 sentence Description]
            | [1 short Insight] | [1 short Recommendation]

Choose a suggestive hex color for each cluster that reflects its clinical
profile (e.g. red tones for high-risk, green for healthy, blue for elderly
chronic, orange for moderate, purple for psychiatric, teal for maternal).
The outlier cluster (-1) should always be gray (#9CA3AF).
\end{lstlisting}

In order for the model to have sufficient context with which to select a color, a title, and a pertinent clinical insight, all features calculated for each cluster in the previous step are concatenated to the prompt in their verbatim form. For each cluster, six lines are added: the important average values (\texttt{age}, \texttt{qaly}, \texttt{daly}, \texttt{total\_procedures}, \texttt{total\_medications}, \texttt{active\_\allowbreak conditions\_\allowbreak count}, \texttt{lifetime\_\allowbreak medication\_\allowbreak days}, \texttt{total\_\allowbreak encounter\_\allowbreak duration\_\allowbreak hours}), the demographics and disease-prevalence dictionary (mode gender, mode marital status, mode encounter class, mode care plan, prevalence of diabetes, hypertension, heart disease, cancer, polypharmacy, BMI distribution, BP distribution, glucose distribution, substance use distribution, readmission distribution), the careplan share, and the top conditions, medications, and procedures pulled from the raw clinical tables joined onto the cluster labels:

\begin{lstlisting}[style=pythonstyle]
prompt = f"""We clustered {len(df_clustered)} patients into groups, plus
{len(df_outliers)} outlier patients. The outliers are labeled as Cluster -1.

Your task is to write extremely brief text to fit inside small frontend UI cards.
... (formatting + colour palette rules omitted) ...

=========================================
CLUSTER DATA SUMMARIES:
=========================================
"""
for c in clusters:
    c_df   = df_all[df_all['cluster'] == c]
    c_size = len(c_df)
    c_pct  = (c_size / total_patients) * 100

    c_means = c_df[['age', 'qaly', 'daly',
                    'total_procedures', 'total_medications',
                    'active_conditions_count', 'lifetime_medication_days',
                    'total_encounter_duration_hours']].mean().round(1).to_dict()

    c_cats = {
        'dominant_gender':    c_df['gender'].mode()[0],
        'dominant_marital':   c_df['marital'].mode()[0],
        'main_encounter_type': c_df['dominant_encounter_class'].mode()[0],
        'dominant_careplan':  c_df['dominant_careplan'].mode()[0],
        'diabetes_rate':      f"{(c_df['has_diabetes']      == 'Yes').mean()*100:.0f}%",
        'hypertension_rate':  f"{(c_df['has_hypertension']  == 'Yes').mean()*100:.0f}%",
        'heart_disease_rate': f"{(c_df['has_heart_disease'] == 'Yes').mean()*100:.0f}%",
        'cancer_rate':        f"{(c_df['has_cancer']        == 'Yes').mean()*100:.0f}%",
        'polypharmacy_rate':  f"{(c_df['is_polypharmacy']   == 'Yes').mean()*100:.0f}%",
        'bmi_mix':            c_df['bmi_type'           ].value_counts(normalize=True).round(2).to_dict(),
        'bp_mix':             c_df['bp_risk_category'   ].value_counts(normalize=True).round(2).to_dict(),
        'glucose_mix':        c_df['glucose_risk_category'].value_counts(normalize=True).round(2).to_dict(),
        'substance_use_mix':  c_df['substance_use_profile'].value_counts(normalize=True).round(2).to_dict(),
        'readmission_mix':    c_df['readmission_risk_tier'].value_counts(normalize=True).round(2).to_dict(),
    }
    careplan_mix = c_df['dominant_careplan'].value_counts(normalize=True).round(2).head(4).to_dict()
    c_cond = top_conditions [top_conditions ['cluster'] == c]['description'].head(4).tolist()
    c_meds = top_medications[top_medications['cluster'] == c]['description'].head(3).tolist()
    c_proc = top_procedures [top_procedures ['cluster'] == c]['description'].head(3).tolist()

    prompt += f"\n### DATA FOR CLUSTER {c} ({c_size} patients, {c_pct:.1f}% of total) ###\n"
    prompt += f"- Key Averages: {c_means}\n"
    prompt += f"- Demographics & Disease Rates: {c_cats}\n"
    prompt += f"- Careplan Mix (share of patients): {careplan_mix}\n"
    prompt += f"- Top Conditions:  {c_cond}\n"
    prompt += f"- Top Medications: {c_meds}\n"
    prompt += f"- Top Procedures:  {c_proc}\n"
\end{lstlisting}

The reason to include the entire dictionary, with all values of \texttt{c\_means}, \texttt{c\_cats}, and \texttt{careplan\_mix} dictionaries in the prompt, is the fact that the LLM must not speculate on anything but the information it gets. The exact numerical measures a clinician would give when describing the cluster, in the same exact units (\texttt{"Obese II"}, \texttt{"Chronic Cardio-Metabolic"}, \texttt{73\%}), are used as input so that the system generates a fitting color and a specific cluster name (red for a cardiometabolic cluster, teal for a maternal one, etc.). The output from the model is then parsed by \texttt{\_parse\_cluster\_summaries} (a simple regex for cluster ID within brackets followed by \texttt{split('|')}), and the result is saved literally into \texttt{cluster} table - that's the exact data you see in cluster summaries on the dashboard.

The same per-cluster values (average age, average BMI, systolic BP, glucose level, QALY/DALY, number of active diagnoses, medications, procedures, days since last encounter, gender balance, smoker distribution, dominant care plan and encounter class, and percentages of diabetes, hypertension, heart disease, cancer, and polypharmacy patients) are requested by the other prompt, the on-demand prompt \texttt{PROMPT\_AI\_ANALYSIS} defined in \texttt{cluster\_service.py}, that will be triggered on the first visit of the detailed cluster page by a user. Compared to the dashboard-card prompt, the novelty here is the inclusion of each metric alongside its population average, allowing the LLM to describe each deviation in relative terms (e.g., "higher/lower than the average"), and the injection of the silhouette and Davies--Bouldin scores available via \texttt{cluster\_meta.metrics} so the model has an idea about the credibility of the cluster. The expected result is again a strict JSON to be processed by \texttt{parse\_ai\_json}, then stored as \texttt{cluster.ai\_analysis} (JSONB):

\begin{lstlisting}[style=pythonstyle]
PROMPT_AI_ANALYSIS = """You are a senior clinical data analyst reviewing a
specific patient group identified through K-Prototypes population segmentation.

IMPORTANT RULES:
- Base every claim strictly on the numbers provided.
- Gender: "balanced" if 40-60%, "more female/male" if 60-70%, "predominantly
  female/male" only if >70%.
- Comparisons: use the "vs population avg" deltas to say above/below average.
- Translate vitals into plain risk language (BP >=140 = hypertensive,
  glucose >=126 = diabetic range, BMI >=30 = obese).

POPULATION CONTEXT:
- {n_patients} patients across {n_clusters} clusters
  | Silhouette: {silhouette} | Davies-Bouldin: {davies_bouldin}
- Population averages: Age {pop_avg_age} | BMI {pop_avg_bmi}
  | Systolic BP {pop_avg_systolic_bp} | Glucose {pop_avg_glucose}
  | QALY {pop_avg_qaly} | DALY {pop_avg_daly}
  | Active conditions {pop_avg_conditions} | Meds {pop_avg_meds}
  | Procedures {pop_avg_procedures}
  | Days since last encounter {pop_avg_days_since}
- Population disease rates: Diabetes {pop_diabetes_pct}%
  | Hypertension {pop_hypertension_pct}% | Heart {pop_heart_pct}%
  | Cancer {pop_cancer_pct}% | Polypharmacy {pop_polypharmacy_pct}%

THIS CLUSTER - "{title}" ({cluster_type}, {count} patients, {pct}% of population):
DEMOGRAPHICS:
- Mean age: {avg_age} (vs pop {pop_avg_age})
- Gender: F {female_pct}% / M {male_pct}%  | Marital mode: {dominant_marital}
VITALS & QUALITY OF LIFE:
- BMI {avg_bmi} (vs {pop_avg_bmi}) | Systolic BP {avg_systolic_bp} (vs {pop_avg_systolic_bp})
- Glucose {avg_glucose} (vs {pop_avg_glucose})
- QALY {avg_qaly} (vs {pop_avg_qaly}) | DALY {avg_daly} (vs {pop_avg_daly})
UTILIZATION & BURDEN:
- Active conditions {avg_conditions} (vs {pop_avg_conditions})
- Medications {avg_meds} (vs {pop_avg_meds})
- Procedures {avg_procedures} (vs {pop_avg_procedures})
- Days since last encounter {avg_days_since} (vs {pop_avg_days_since})
- Dominant encounter class: {dominant_encounter}
- Dominant careplan: {dominant_careplan}
DISEASE PREVALENCE:
- Diabetes {diabetes_pct}% (vs pop {pop_diabetes_pct}%)
- Hypertension {hypertension_pct}% (vs pop {pop_hypertension_pct}%)
- Heart disease {heart_disease_pct}% (vs pop {pop_heart_pct}%)
- Cancer {cancer_pct}% (vs pop {pop_cancer_pct}%)
- Polypharmacy {polypharmacy_pct}% (vs pop {pop_polypharmacy_pct}%)
- Smokers: current {current_smoker_pct}% / former {former_smoker_pct}%
  / never {never_smoker_pct}%

TASK: Analyse this patient group for TWO audiences:
1. CLINICIANS - 2-3 sentences anchored on dominant careplan + top prevalence
   deltas, the main risks, and the most actionable care gaps.
2. GENERAL PUBLIC - 1-2 sentences in plain everyday language, no medical terms.

Respond ONLY with valid JSON:
{{
  "risk_level": "high or moderate or low",
  "medical_interpretation": "...",
  "layperson_summary": "...",
  "key_characteristics":   ["...", "...", "...", "..."],
  "priority_interventions": ["...", "...", "..."],
  "care_recommendations":  ["...", "...", "..."]
}}
"""
\end{lstlisting}

\noindent The combination of both prompts makes up the AI part of the clusters dashboard; the first prompt is executed once per training to generate the name + color + one-liner that will be displayed on each individual card, whereas the second prompt is generated lazily after each user click to generate the actual medical interpretation, the layman’s version, and the action items list.

\subsubsection*{Libraries}

The principal libraries used are \texttt{kmodes} (algorithm), \texttt{scikit-learn} (\texttt{StandardScaler}, \texttt{MinMaxScaler}, \texttt{RobustScaler}, \texttt{PCA}, evaluation metrics), \texttt{prince} (\texttt{FAMD}), \texttt{matplotlib} and \texttt{seaborn} (plots), \texttt{pandas}/\texttt{numpy} (feature pipeline), \texttt{SQLAlchemy} (read/write paths), and the OpenAI Python SDK behind \texttt{ask\_ai} for the cluster-narrative step.

